\import{set/equinumerosity-product.tex}
\import{set/equinumerosity-union.tex}
\import{set/suc.tex}

\section{Tagged disjoint sums}

\begin{definition}\label{sum_left_tag}
    $\sumLeftTag=\emptyset$.
\end{definition}

\begin{definition}\label{sum_right_tag}
    $\sumRightTag=\suc{\emptyset}$.
\end{definition}

\begin{definition}\label{tagged_sum}
    $\taggedSum{A}{B}
        =(A\times\{\sumLeftTag\})\union
            (B\times\{\sumRightTag\})$.
\end{definition}

\begin{proposition}\label{sum_tags_distinct}
    $\sumLeftTag\neq\sumRightTag$.
\end{proposition}
\begin{proof}
    Follows by \obv{sum_left_tag,sum_right_tag,suc_neq_emptyset}.
\end{proof}

\begin{proposition}\label{tagged_sum_components_disjoint}
    $(A\times\{\sumLeftTag\})\inter
        (B\times\{\sumRightTag\})=\emptyset$.
\end{proposition}
\begin{proof}
    Show that every element of
        $(A\times\{\sumLeftTag\})\inter
            (B\times\{\sumRightTag\})$ is an element of $\emptyset$.
    \begin{subproof}
        Fix $p\in
            (A\times\{\sumLeftTag\})\inter
                (B\times\{\sumRightTag\})$.
        Hence $p\in A\times\{\sumLeftTag\}$ and
            $p\in B\times\{\sumRightTag\}$ by \obv{inter}.
        Hence $\snd{p}=\sumLeftTag$ and $\snd{p}=\sumRightTag$
            by \auto{snd_type,singleton_iff}.
        Contradiction by \obv{sum_tags_distinct}.
    \end{subproof}
    Hence
        $(A\times\{\sumLeftTag\})\inter
            (B\times\{\sumRightTag\})\subseteq\emptyset$
        by \obv{subseteq}.
    Follows by \obv{subseteq_emptyset_iff}.
\end{proof}

\begin{theorem}\label{tagged_sum_equinum_congruent}
    Suppose $A\approx C$ and $B\approx D$.
    Then $\taggedSum{A}{B}\approx\taggedSum{C}{D}$.
\end{theorem}
\begin{proof}
    $\{\sumLeftTag\}\approx\{\sumLeftTag\}$ and
        $\{\sumRightTag\}\approx\{\sumRightTag\}$
        by \obv{equinum_refl}.
    Hence $A\times\{\sumLeftTag\}
            \approx C\times\{\sumLeftTag\}$ and
        $B\times\{\sumRightTag\}
            \approx D\times\{\sumRightTag\}$
        by \obv{times_equinum_congruent}.
    $(A\times\{\sumLeftTag\})\inter
        (B\times\{\sumRightTag\})=\emptyset$
        by \obv{tagged_sum_components_disjoint}.
    $(C\times\{\sumLeftTag\})\inter
        (D\times\{\sumRightTag\})=\emptyset$
        by \obv{tagged_sum_components_disjoint}.
    Follows by
        \auto{tagged_sum,disjoint_union_equinum_congruent}.
\end{proof}

\begin{theorem}\label{tagged_sum_injection_congruent}
    Let $f\in\inj{A}{C}$.
    Let $g\in\inj{B}{D}$.
    Then there exists $h\in\inj{\taggedSum{A}{B}}{\taggedSum{C}{D}}$.
\end{theorem}
\begin{proof}
    $\identity{\{\sumLeftTag\}}
        \in\inj{\{\sumLeftTag\}}{\{\sumLeftTag\}}$
        by
        \auto{id_elem_funs,identity_injective,inj,funs_is_function,injective_function}.
    Take $f_L$ such that $f_L\in\inj{A\times\{\sumLeftTag\}}
        {C\times\{\sumLeftTag\}}$
        by \auto{times_injection_congruent}.
    $\identity{\{\sumRightTag\}}
        \in\inj{\{\sumRightTag\}}{\{\sumRightTag\}}$
        by
        \auto{id_elem_funs,identity_injective,inj,funs_is_function,injective_function}.
    Take $g_R$ such that $g_R\in\inj{B\times\{\sumRightTag\}}
        {D\times\{\sumRightTag\}}$
        by \auto{times_injection_congruent}.
    $f_L\in\funs{A\times\{\sumLeftTag\}}
        {C\times\{\sumLeftTag\}}$ by \obv{inj}.
    $g_R\in\funs{B\times\{\sumRightTag\}}
        {D\times\{\sumRightTag\}}$ by \obv{inj}.
    $C\times\{\sumLeftTag\}\subseteq\taggedSum{C}{D}$
        by \auto{tagged_sum,union_upper_left}.
    Hence $f_L\in\inj{A\times\{\sumLeftTag\}}{\taggedSum{C}{D}}$
        by \obv{inj_weaken_codom}.
    $D\times\{\sumRightTag\}\subseteq\taggedSum{C}{D}$
        by \auto{tagged_sum,union_upper_right}.
    Hence $g_R\in\inj{B\times\{\sumRightTag\}}{\taggedSum{C}{D}}$
        by \obv{inj_weaken_codom}.
    $\ran{f_L}\subseteq C\times\{\sumLeftTag\}$ and
        $\ran{g_R}\subseteq D\times\{\sumRightTag\}$
        by \obv{inj,funs_ran}.
    Show that every element of $\ran{f_L}\inter\ran{g_R}$
        is an element of $\emptyset$.
    \begin{subproof}
        Fix $z\in\ran{f_L}\inter\ran{g_R}$.
        Hence $z\in\ran{f_L}$ and $z\in\ran{g_R}$ by \obv{inter}.
        Hence $z\in C\times\{\sumLeftTag\}$ and
            $z\in D\times\{\sumRightTag\}$ by \obv{subseteq}.
        Hence $z\in
            (C\times\{\sumLeftTag\})\inter
                (D\times\{\sumRightTag\})$ by \obv{inter_intro}.
        Hence $z\in\emptyset$
            by \obv{tagged_sum_components_disjoint}.
    \end{subproof}
    Hence $\ran{f_L}\inter\ran{g_R}\subseteq\emptyset$
        by \obv{subseteq}.
    Hence $\ran{f_L}\inter\ran{g_R}=\emptyset$
        by \obv{subseteq_emptyset_iff}.
    $(A\times\{\sumLeftTag\})\inter
        (B\times\{\sumRightTag\})=\emptyset$
        by \obv{tagged_sum_components_disjoint}.
    Hence $f_L\union g_R
        \in\inj{\taggedSum{A}{B}}{\taggedSum{C}{D}}$
        by \auto{tagged_sum,union_inj_disjoint}.
\end{proof}

\begin{proposition}\label{tagged_sum_equinum_commutative}
    $\taggedSum{A}{B}\approx\taggedSum{B}{A}$.
\end{proposition}
\begin{proof}
    $A\times\{\sumLeftTag\}\approx A$
        by \obv{times_singleton_right_equinum}.
    $A\approx A\times\{\sumRightTag\}$
        by \obv{times_singleton_right_equinum,equinum_sym}.
    Hence $A\times\{\sumLeftTag\}
        \approx A\times\{\sumRightTag\}$ by \obv{equinum_tran}.
    $B\times\{\sumRightTag\}\approx B$
        by \obv{times_singleton_right_equinum}.
    $B\approx B\times\{\sumLeftTag\}$
        by \obv{times_singleton_right_equinum,equinum_sym}.
    Hence $B\times\{\sumRightTag\}
        \approx B\times\{\sumLeftTag\}$ by \obv{equinum_tran}.
    $(A\times\{\sumLeftTag\})\inter
        (B\times\{\sumRightTag\})=\emptyset$
        by \obv{tagged_sum_components_disjoint}.
    $(A\times\{\sumRightTag\})\inter
        (B\times\{\sumLeftTag\})=\emptyset$
        by \auto{inter_comm,tagged_sum_components_disjoint}.
    Hence
        $(A\times\{\sumLeftTag\})\union
            (B\times\{\sumRightTag\})
        \approx
        (A\times\{\sumRightTag\})\union
            (B\times\{\sumLeftTag\})$
        by \obv{disjoint_union_equinum_congruent}.
    $(A\times\{\sumRightTag\})\union
        (B\times\{\sumLeftTag\})
        =(B\times\{\sumLeftTag\})\union
            (A\times\{\sumRightTag\})$
        by \obv{union_comm}.
    Follows by \auto{tagged_sum}.
\end{proof}

\begin{proposition}\label{tagged_sum_empty_right_equinum}
    $\taggedSum{A}{\emptyset}\approx A$.
\end{proposition}
\begin{proof}
    $\emptyset\times\{\sumRightTag\}=\emptyset$
        by \obv{times_empty_left}.
    Hence $\taggedSum{A}{\emptyset}
        =A\times\{\sumLeftTag\}$
        by \auto{tagged_sum,union_emptyset}.
    Follows by \auto{times_singleton_right_equinum}.
\end{proof}

\begin{corollary}\label{tagged_sum_empty_left_equinum}
    $\taggedSum{\emptyset}{A}\approx A$.
\end{corollary}
\begin{proof}
    $\taggedSum{\emptyset}{A}\approx\taggedSum{A}{\emptyset}$
        by \obv{tagged_sum_equinum_commutative}.
    $\taggedSum{A}{\emptyset}\approx A$
        by \obv{tagged_sum_empty_right_equinum}.
    Follows by \obv{equinum_tran}.
\end{proof}

\begin{proposition}\label{tagged_sum_equinum_associative}
    $\taggedSum{\taggedSum{A}{B}}{C}
        \approx\taggedSum{A}{\taggedSum{B}{C}}$.
\end{proposition}
\begin{proof}
    Let $L_A=(A\times\{\sumLeftTag\})\times\{\sumLeftTag\}$.
    Let $L_B=(B\times\{\sumRightTag\})\times\{\sumLeftTag\}$.
    Let $L_C=C\times\{\sumRightTag\}$.
    Let $R_A=A\times\{\sumLeftTag\}$.
    Let $R_B=(B\times\{\sumLeftTag\})\times\{\sumRightTag\}$.
    Let $R_C=(C\times\{\sumRightTag\})\times\{\sumRightTag\}$.
    $\taggedSum{\taggedSum{A}{B}}{C}
        =(L_A\union L_B)\union L_C$
        by \auto{tagged_sum,union_times_left}.
    $\taggedSum{A}{\taggedSum{B}{C}}
        =(R_A\union R_B)\union R_C$
        by \auto{tagged_sum,union_times_left,union_assoc}.

    $L_A\approx R_A$ by \auto{times_singleton_right_equinum}.
    $L_B\approx B\times\{\sumRightTag\}$
        by \auto{times_singleton_right_equinum}.
    $B\times\{\sumRightTag\}\approx B$
        by \obv{times_singleton_right_equinum}.
    $R_B\approx B\times\{\sumLeftTag\}$
        by \auto{times_singleton_right_equinum}.
    $B\times\{\sumLeftTag\}\approx B$
        by \obv{times_singleton_right_equinum}.
    Hence $B\approx R_B$
        by \obv{equinum_sym,equinum_tran}.
    Hence $L_B\approx R_B$ by \obv{equinum_tran}.
    $R_C\approx L_C$ by \auto{times_singleton_right_equinum}.
    Hence $L_C\approx R_C$ by \obv{equinum_sym}.

    $L_A\inter L_B=\emptyset$
        by
        \auto{inter_times,tagged_sum_components_disjoint,times_empty_left}.
    $L_A\inter L_C=\emptyset$
        by \auto{tagged_sum_components_disjoint}.
    $L_B\inter L_C=\emptyset$
        by \auto{tagged_sum_components_disjoint}.
    $R_A\inter R_B=\emptyset$
        by \auto{tagged_sum_components_disjoint}.
    $R_A\inter R_C=\emptyset$
        by \auto{tagged_sum_components_disjoint}.
    $R_B\inter R_C=\emptyset$
        by
        \auto{inter_times,tagged_sum_components_disjoint,times_empty_left}.
    Hence $(L_A\union L_B)\union L_C
        \approx(R_A\union R_B)\union R_C$
        by \obv{three_disjoint_unions_equinum_congruent}.
\end{proof}

\begin{proposition}\label{times_tagged_sum_equinum_distributive}
    $A\times\taggedSum{B}{C}
        \approx\taggedSum{A\times B}{A\times C}$.
\end{proposition}
\begin{proof}
    Let $S_B=A\times(B\times\{\sumLeftTag\})$.
    Let $S_C=A\times(C\times\{\sumRightTag\})$.
    Let $T_B=(A\times B)\times\{\sumLeftTag\}$.
    Let $T_C=(A\times C)\times\{\sumRightTag\}$.
    $A\times\taggedSum{B}{C}=S_B\union S_C$
        by \auto{tagged_sum,union_times_right}.
    $\taggedSum{A\times B}{A\times C}=T_B\union T_C$
        by \obv{tagged_sum}.
    $T_B\approx S_B$ by \auto{times_equinum_associative}.
    Hence $S_B\approx T_B$ by \obv{equinum_sym}.
    $T_C\approx S_C$ by \auto{times_equinum_associative}.
    Hence $S_C\approx T_C$ by \obv{equinum_sym}.
    $S_B\inter S_C=\emptyset$
        by
        \auto{inter_times,inter_idempotent,tagged_sum_components_disjoint,times_empty_right}.
    $T_B\inter T_C=\emptyset$
        by \auto{tagged_sum_components_disjoint}.
    Hence $S_B\union S_C\approx T_B\union T_C$
        by \obv{disjoint_union_equinum_congruent}.
\end{proof}
